\SourceLeaf{198}
\Needspace{8\baselineskip}
\pdfbookmark[1]{Seeber, positive ternäre quadratische Formen. Anzeige 1831 Juli 9}{notice-1831-07-09}
\begin{center}
\rule{.86\linewidth}{.3pt}\par
Göttingische gelehrte Anzeigen.\quad 1831 Juli 9.\par
\rule{.86\linewidth}{.3pt}
\end{center}
\emph{Untersuchungen über die Eigenschaften der positiven ternären quadratischen Formen von} \textsc{Ludwig August Seeber}, \emph{Dr. der Philosophie, ordentl. Professor der Physik an der Universität in Freiburg. Freiburg im Breisgau} 1831. (248 S. in 4.)

Die Functionen zweier unbestimmten Grössen $x$ und $y$ von der Gestalt $axx+2bxy+cyy$, wo $a,b,c$ bestimmte ganze Zahlen vorstellen, bilden bekanntlich unter dem Namen der \emph{quadratischen Formen}, oder, wo eine weitere Unterscheidung erforderlich wird, der \emph{binären quadratischen Formen}, einen der interessantesten und reichhaltigsten Gegenstände der höheren Arithmetik. Die dabei zunächst vorkommenden Aufgaben: zu entscheiden, ob eine solche gegebene Form eine andere $a'x'x'+2b'x'y'+c'y'y'$ unter sich begreift, d. i. durch eine Substitution $x=\alpha x'+\beta y'$, $y=\gamma x'+\delta y'$, in welcher $\alpha,\beta,\gamma,\delta$ ganze Zahlen sind, in dieselbe verwandelt werden kann; ob eine solche Relation zweier Formen eine gegenseitige ist, wo die Formen äquivalent heissen; ferner in beiden Fällen alle möglichen Umformungen der einen in die andere anzugeben; endlich alle möglichen Darstellungen einer gegebenen ganzen Zahl durch eine gegebene Form vermöge ganzer Werthe der unbestimmten Grössen aufzufinden --- diese Aufgaben sind in den \emph{Disquisitiones Arithmeticae} vollständig aufgelöset, machen aber von dem die quadratischen Formen betreffenden Abschnitte dieses Werks nur den bei weitem kleineren Theil aus. Die darauf folgenden feineren Untersuchungen erfordern\textsuperscript{.} zum Theil eine vorläufige Bearbeitung eines um eine Stufe höheren und viel grössere Schwierigkeiten darbietenden Feldes, nemlich der Lehre von ähnlichen Functionen dreier unbestimmter Grössen $x,y,z$, welche also die Gestalt haben $axx+byy+czz+2a'yz+2b'xz+2c'xy$, und ternäre quadratische Formen heissen. Die Auflösung der diese ternären Formen betreffenden Hauptaufgaben ist in dem erwähnten Werke entwickelt, jedoch nur so weit, als zu dem angezeigten Zwecke nothwendig war. Nach einem Zwischenraum von dreissig Jahren hat nun der Verfasser des vorliegenden Werks zuerst diese Untersuchungen wieder aufgenommen, und in Beziehung auf die eine Hauptgattung der ternären Formen, nemlich die positiven, dasjenige was in den \emph{Disquisitiones }
\SourceLeaf{199}
\emph{Arithmeticae} unvollendet gelassen war, zur Vollständigkeit gebracht. Für diejenigen, welche aus der höheren Arithmetik ein tieferes Studium gemacht haben, würden wir dasjenige, was in dem vorliegenden Werke Neues geleistet ist, mit wenigen Worten bezeichnen können; allein, um auch andern verständlich zu sein, müssen wir uns etwas mehr Ausführlichkeit verstatten, und wir thun dies um so lieber, da diese Untersuchungen auch ausserhalb des Gebietes der höheren Arithmetik ein eigenthümliches Interesse haben.

Die Eigenschaften einer binären Form $axx+2bxy+cyy$ hängen vornehmlich von der Zahl $bb-ac$ ab, welche daher der Determinant jener Form heisst. Zwei äquivalente Formen haben allemal gleiche Determinanten. Allein nicht alle Formen, die einen gegebenen Determinanten haben, sind darum schon äquivalent, vielmehr zerfallen solche Formen in eine kleinere oder grössere, aber stets endliche Anzahl von Classen, so dass die zu einerlei Classe gehörigen unter sich äquivalent, die zu verschiedenen Classen gehörenden hingegen nicht äquivalent sind. Durch Formen, deren Determinant positiv ist, lassen sich ohne Unterschied positive und negative Zahlen darstellen; hingegen durch Formen mit negativem Determinanten sind nur solche Zahlen darstellbar, welche mit $a$ und $c$ einerlei Zeichen haben, daher hier positive und negative Formen unterschieden werden. Die einfachsten Formen in jeder Classe haben bestimmte Kriterien, heissen reducirte Formen, und können als Repräsentanten der ganzen Classe betrachtet werden.

Aehnliche Verhältnisse in Beziehung auf die ternären Formen sind in den \emph{Disquisitiones Arithmeticae} nachgewiesen. Determinant der ternären Form
\[
axx+byy+czz+2a'yz+2b'xz+2c'xy
\]
heisst die Zahl
\[
aa'a'+bb'b'+cc'c'-abc-2a'b'c'
\]
Auch hier ist zur Aequivalenz zweier Formen die Gleichheit der Determinanten erforderlich, aber nicht zureichend, sondern sämmtliche Formen mit einem bestimmten Determinanten zerfallen in eine endliche Anzahl von Classen, in deren jeder die einfachsten Formen reducirte heissen können und alle übrigen gleichsam repräsentiren. Mit dem Unterschiede zwischen positiven und negativen Formen verhält es sich aber hier anders, als bei den binären Formen. Für jeden gegebenen Determinanten, er sei positiv oder negativ, gibt es theils Formen, durch welche
\SourceLeaf{200}
ohne Unterschied positive und negative Zahlen darstellbar sind (indifferente Formen), theils solche Formen, durch die entweder nur positive oder nur negative Zahlen sich darstellen lassen (positive oder negative Formen); allein positive Formen gibt es nur für negative Determinanten, und negative nur für positive. Uebrigens ist es von selbst klar, dass die Qualification einer Form, insofern sie indifferent, positiv oder negativ ist, zugleich der ganzen Classe, zu welcher sie gehört, zukommt. Das vorliegende Werk beschränkt sich auf die positiven Formen, deren Determinanten also negativ sein müssen: offenbar findet aber alles, was von diesen gilt, von selbst seine Uebertragung auf die negativen Formen, während die in dem Werke ganz ausgeschlossenen indifferenten Formen eine ganz abweichende Behandlung erfordern.

In den \emph{Disquisitiones Arithmeticae} war, wie schon erwähnt ist, die Theorie der ternären Formen nur so weit entwickelt, als für den dortigen Zweck nöthig war, und daher die Aufgabe, die Aequivalenz zweier gegebenen ternären Formen zu entscheiden, noch nicht in vollständiger Allgemeinheit aufgelöset. Zwar war daselbst gezeigt, wie man zu jeder vorgegebenen Form eine äquivalente der einfachsten Art finden, und dass es solcher reducirten Formen für jeden gegebenen Determinanten nur eine endliche Anzahl geben könne; allein da es in jeder Classe mehrere solcher reducirten Formen gibt, die sich nicht in allen Fällen \emph{sogleich} als äquivalent ergeben, so fehlte noch ein Kriterium, woran man die Aequivalenz oder Nicht-Aequivalenz solcher Formen mit Gewissheit erkennen kann. Dieses Bedürfniss hat nun der Verfasser des vorliegenden Werks in Beziehung auf die positiven Formen vollständig und mit musterhafter Gründlichkeit gehoben. Sein Verfahren ist übrigens etwas anders eingekleidet, als wir die Sache so eben ausgesprochen haben, und wie sie sich verhalten müsste, wenn man in den Begriff der reducirten positiven Formen nur die wesentlichsten Bedingungen der grössten Einfachheit aufnimmt, welche in dem Fall der positiven Formen die sind, dass die (ihrer Natur nach positiven) Zahlen $a,b,c$ nicht kleiner sein dürfen, als respective $b'$ oder $c'$, $a'$ oder $c'$, $a'$ oder $b'$ ohne Rücksicht auf die Zeichen. Herr \textsc{Seeber} hat nemlich dem Begriff der reducirten Formen noch solche Modificationen hinzugesetzt, dass es in jeder Classe immer nur Eine der Art geben kann, Eine aber geben muss. Wegen eines schönen von Herrn \textsc{Seeber} durch Induction gefundenen weiter unten noch zu erwähnenden Theorems führen wir hier die Hauptbedingungen, welche Hr. S. in den Begriff der reducirten Formen 
\SourceLeaf{201}
aufgenommen hat, an: diese sind 1) dass unter den Zahlen $a',b',c'$ nicht zwei von entgegengesetzten Zeichen sein dürfen; 2) dass ohne Rücksicht auf das Zeichen $2b'$ und $2c'$ nicht grösser als $a$ sein dürfen, ferner $a$ und $2a'$ nicht grösser als $b$, und $b$ nicht grösser als $c$; 3) dass in dem Fall, wo $a',b',c'$ zugleich negativ sind, die doppelte Summe dieser Zahlen nicht grösser als $a+b$ sein darf. Die übrigen noch für einige specielle Fälle hinzukommenden Modificationen können wir hier übergehen.

Den Hauptinhalt des Werkes macht nun zuerst die Auflösung der Aufgabe aus, zu jeder gegebenen positiven Form eine äquivalente zu finden, die nach der festgesetzten Definition den Character einer reducirten hat, und dann der strenge Beweis des Lehrsatzes, dass zwei nicht identische reducirte Formen nicht äquivalent sein können, oder was dasselbe ist, dass es in jeder Classe nur eine reducirte Form gibt. Dem Geiste der Gründlichkeit, womit diese Gegenstände durchgeführt sind, müssen wir volle Gerechtigkeit widerfahren lassen, und wenn wir es dabei bedauern müssen, dass damit eine sehr grosse und vielleicht manchen abschreckende Weitläuftigkeit verbunden gewesen ist, da die Auflösung des Problems 41 Seiten, und der Beweis des Theorems 91 Seiten einnimmt, so wollen wir dies doch keinesweges als einen Tadel angesehen wissen. Wenn ein schwieriges Problem oder Theorem aufzulösen oder zu beweisen vorliegt, so ist allezeit der erste und mit gebührendem Danke zu erkennende Schritt, dass überhaupt eine Auflösung oder ein Beweis gefunden werde, und die Frage, ob dies nicht auf eine leichtere und einfachere Art hätte geschehen können, bleibt so lange eine müssige, als die Möglichkeit nicht zugleich durch die That entschieden wird. Wir halten es daher für unzeitig, hier bei dieser Frage zu verweilen. --- Der übrige Theil des Werkes enthält noch hauptsächlich die mit gleicher Gründlichkeit durchgeführten Auflösungen der Aufgaben: zu entscheiden, ob eine gegebene Form eine andere gegebene ihr nicht äquivalente unter sich begreife; alle möglichen Transformationen einer gegebenen Form in eine gegebene äquivalente oder nur unter ihr begriffene zu finden; endlich für einen gegebenen Determinanten alle möglichen Classen positiver ternärer Formen anzugeben.

Wir müssen noch bemerken, dass Herr \textsc{Seeber} die Gestalt der ternären Formen etwas anders gefasst hat, als in den \emph{Disquisitiones Arithmeticae} geschehen war, wo, mit Vorbedacht, die Coëfficienten der Producte $yz$, $xz$, $xy$ als gerade Zahlen vorausgesetzt waren, wogegen Hr. S. auch ungerade zulässt, und daher
\SourceLeaf{202}
mit $a',b',c'$ bezeichnet, was oben mit $2a'$, $2b'$, $2c'$ bezeichnet war. Offenbar ist die grössere Allgemeinheit, welche dadurch erreicht wird, nur scheinbar, oder doch überflüssig, da alles was von solchen Formen mit ungeraden Coëfficienten gesagt werden kann, sich auch von selbst ergibt, wenn man anstatt derselben ihr Doppeltes in Betracht zieht: wir können daher diese Abänderung, wodurch überdies einiger Verlust an Einfachheit entsteht, nicht billigen. Eine Folge davon ist gewesen, dass das, was Herr \textsc{Seeber} Determinant nennt, allemal das Vierfache von der Zahl ist, welche in den \emph{Disquisitiones Arithmeticae} diesen Namen führt. In gegenwärtiger Anzeige haben wir die Terminologie der \emph{Disquisitiones Arithmeticae} beibehalten.

Bei dem zuletzt erwähnten Problem (zu jedem gegebenen Determinanten alle möglichen reducirten Formen anzugeben) hat Herr \textsc{Seeber}, um Grenzen für die drei ersten Coëfficienten zu haben, ein Theorem benutzt, vermöge dessen das Product derselben $abc$ nicht grösser sein kann, als der dreifache Determinant. Dieses Theorem ist von Hn. \textsc{Seeber} strenge bewiesen; allein in der Vorrede bemerkt er, dass er unter mehr als 600 von ihm untersuchten Fällen nicht einen einzigen gefunden habe, wo jenes Product das Doppelte des Determinanten überschritten hätte, und hält es daher für höchst wahrscheinlich, dass diese engere Begrenzung allgemeingültig sei; es sei ihm jedoch nicht gelungen, einen strengen Beweis dafür zu finden. Da dieses auf dem Wege der Induction von Herrn \textsc{Seeber} gefundene Theorem sowohl an sich merkwürdig, als für die Abkürzung der Auflösung der erwähnten Aufgabe wichtig ist, so wollen wir hier, um auch unsererseits in dieser Anzeige einen Beitrag zur Vervollkommnung dieser Theorie zu geben, einen sehr einfachen Beweis beifügen. Es müssen dabei zwei Fälle unterschieden werden.

I. Wenn von den Zahlen $a',b',c'$ keine negativ ist, so setze man
\[
\begin{aligned}
b-2a'&=d,&c-2b'&=e,&a-2c'&=f\\
c-2a'&=g,&a-2b'&=h,&b-2c'&=i
\end{aligned}
\]
wo aus der Definition der reducirten positiven Formen sogleich folgt, dass wenn
\[
axx+byy+czz+2a'yz+2b'xz+2c'xy
\]
eine solche ist, keine jener sechs Zahlen negativ ist, so wie sich von selbst versteht, dass $a,b,c$ positiv sind. Bezeichnet man nun den (negativen) 
\SourceLeaf{203}
Determinanten der Form durch $-D$, so hat man, wie man sich durch die Entwickelung leicht überzeugt, die identische Gleichung
\[
2D-abc=aa'd+bb'e+cc'f+a'hi+b'gi+c'gh+ghi
\]
in welcher keines der sieben Glieder zur Rechten negativ sein kann, und folglich $abc$ nicht grösser als $2D$. Dasselbe folgt auf gleiche Weise aus der identischen Gleichung
\[
2D-abc=aa'g+bb'h+cc'i+a'ef+b'df+c'de+def
\]

II. Wenn keine der Zahlen $a',b',c'$ positiv ist, setze man
\[
\begin{aligned}
b+2a'&=d,&c+2b'&=e,&a+2c'&=f\\
c+2a'&=g,&a+2b'&=h,&b+2c'&=i
\end{aligned}
\]
\[
\begin{aligned}
b+c+2a'+2b'+2c'&=k\\
a+c+2a'+2b'+2c'&=l\\
a+b+2a'+2b'+2c'&=m
\end{aligned}
\]
und den Determinanten der Form wie vorhin $=-D$. Vermöge der Definition der reducirten positiven Formen wird keine der neun Zahlen $d,\allowbreak e,\allowbreak f,\allowbreak g,\allowbreak h,\allowbreak i,\allowbreak k,\allowbreak l,\allowbreak m$ negativ sein können, und so ergibt sich aus der identischen Gleichung
\[
\begin{aligned}
6D-3abc={}&-aa'(d+2k)-bb'(e+2l)-cc'(f+2m)\\
&-a'hi-b'gi-c'gh+def+2ghi
\end{aligned}
\]
in welcher, weil $a',b',c'$ nicht positiv, sondern negativ oder Null sind, alle Glieder zur Rechten positiv oder Null werden, dass $3abc$ nicht grösser als $6D$, oder $abc$ nicht grösser als $2D$ sein kann. Dasselbe folgt eben so aus der identischen Gleichung
\[
\begin{aligned}
6D-3abc={}&-aa'(g+2k)-bb'(h+2l)-cc'(i+2m)\\
&-a'ef-b'df-c'de+2def+ghi
\end{aligned}
\]
Beide Gleichungen sind symmetrisch. Verzichtet man auf völlige Symmetrie, so ist der Beweis mit einer noch geringern Anzahl von Gliedern zu führen, z.~B. durch die identische Gleichung
\[
\begin{aligned}
8D-4abc={}&-2aa'(g+k)-2bb'(e+l)-4cc'm\\
&+(c+e)df+(c+g)hi
\end{aligned}
\]

\SourceLeaf{204}
Wir wollen nun noch einiges über die Bedeutung der positiven binären und ternären quadratischen Formen ausser dem Gebiete der höheren Arithmetik hinzusetzen: von den negativen besonders zu handeln ist unnöthig, und die indifferenten entziehen sich dieser Behandlung ganz.

Die positive binäre Form $axx+2bxy+cyy$ stellt allgemein das Quadrat der Entfernung zweier unbestimmter Punkte in einer Ebene vor, deren Coordinaten in Beziehung auf zwei unter einem Winkel, dessen Cosinus $=\frac{b}{\surd ac}$ ist, gegen einander geneigte Axen um $x\surd a$, $y\surd c$ verschieden sind. Insofern $x$ und $y$ also nur \emph{ganze} Zahlen bedeuten sollen, bezieht sich die Form auf ein System parallelogrammatisch geordneter Punkte, die in den Durchschnitten zweier Systeme von Parallellinien liegen. Die Linien jedes Systems sind in gleichen Entfernungen von einander, und zwar sind die des einen, wenn sie parallel mit den Linien des zweiten gemessen werden, $=\surd a$; die Entfernungen des andern, parallel mit den Linien des ersten gemessen, $=\surd c$: die Neigung beider Systeme gegen einander die oben angegebene. Auf diese Weise erscheint die Ebene in lauter gleiche Parallelogramme getheilt, deren Eckpunkte das Punktensystem ausmachen, ohne dass irgend einer der Punkte innerhalb eines Parallelogrammes fallen kann. Der Determinant mit positivem Zeichen genommen, also $ac-bb$, bedeutet das Quadrat des Flächeninhalts eines Elementar-Parallelogramms. Ein und dasselbe System solcher Punkte kann auf unendlich viele verschiedene Arten parallelogrammatisch abgetheilt, und also auf ebenso viele verschiedene Formen zurückgeführt werden: alle diese verschiedenen Formen sind aber, was in der Kunstsprache äquivalent heisst, und der Inhalt eines Elementar-Parallelogramms bleibt allemal derselbe. Zwei Formen, die nicht äquivalent sind, von denen aber die eine die andere unter sich begreift, beziehen sich auf dasselbe System von Punkten, aber die erstere Form auf das ganze System, die zweite auf einen Theil. Zwei Formen, die, nach der Kunstsprache, uneigentlich äquivalent (improprie aequivalentes) heissen, beziehen sich auf zwei gleiche aber verkehrt liegende Systeme von Punkten, indem man sich die Ebene umgekehrt gelegt denkt u.~s.~w.

Auf gleiche Weise bedeutet allgemein die positive ternäre Form
\[
axx+byy+czz+2a'yz+2b'xz+2c'xy
\]
das Quadrat der Entfernung zweier unbestimmter Punkte im Raume, deren Coordinaten in Beziehung auf drei Axen (1), (2), (3) die Unterschiede $x\surd a$, $y\surd b$, $z\surd c$
\SourceLeaf{205}
geben: die Cosinus der Winkel zwischen den Axen (2) und (3), (1) und (3), (1) und (2) sind hier resp. $\frac{a'}{\surd bc}$, $\frac{b'}{\surd ac}$, $\frac{c'}{\surd ab}$. Insofern hier $x,y,z$ blos ganze Zahlen bedeuten sollen, bezieht sich die Form auf ein System parallelepipedisch geordneter, d.~i. durch die Durchschnitte dreier Systeme paralleler äquidistanter Ebenen sich ergebender Punkte. Der ganze Raum erscheint so in lauter gleiche Parallelepipeden getheilt, deren Eckpunkte jenes System von Punkten ausmachen, und das Quadrat des Rauminhalts eines Elementar-Parallelepipedum ist dem mit positivem Zeichen genommenen\textsuperscript{.} Determinanten der ternären Form gleich. Aequivalente Formen repräsentiren ein und dasselbe System von Punkten, nur auf andere Axen oder Fundamentalebenen bezogen. Auf gleiche Weise finden alle andere Hauptmomente der Theorie der ternären Formen hier ihre geometrische Bedeutung, das Enthaltensein einer Form unter einer andern, die Darstellung einer bestimmten Zahl oder einer unbestimmten binären Form durch eine ternäre, die Lehre von den zugeordneten ternären Formen (formae adiunctae), das Wegfallen der Unterscheidung zwischen eigentlicher und uneigentlicher Aequivalenz, das Wesen der reducirten Formen u.~s.~w., wir müssen uns aber auf obige Andeutungen beschränken, zumal da das vorliegende Werk, welches die ternären Formen lediglich aus rein arithmetischem Gesichtspunkte betrachtet, nur mittelbarer Weise Veranlassung dazu gegeben hat. Man wird wenigstens daraus erkennen, welch ein reiches Feld hier den Untersuchungen geöffnet ist, die nicht blos für sich ein hohes theoretisches Interesse haben, sondern auch zu einer eben so bequemen als allgemeinen Behandlung aller Relationen unter den Krystallformen benutzt werden können. In das Detail dieser Benutzung einzugehen, ist hier der Ort nicht: wir dürfen jedoch die Bemerkung nicht übergehen, dass wenn gleich ursprünglich angenommen ist, dass $a,b,c,a',b',c'$ ganze Zahlen vorstellen, doch der grösste Theil der Lehre von den ternären Formen, und namentlich dasjenige, was für jene Benutzung erforderlich ist, auch unabhängig von jener Voraussetzung gültig bleibt. In der That führen zwar \textsc{Haüy}'s Angaben bei den meisten Krystallgattungen auf sehr einfache ganze Werthe der Coëfficienten in den ternären Formen, welche sich auf die jenen entsprechende Anordnung des Punktensystems beziehen; allein die genaueren spätern Messungen von \textsc{Wollaston}, \textsc{Malus}, \textsc{Biot}, \textsc{Kupffer} u.~a. stehen damit im Widerspruch, und machen es zweifelhaft, ob rationale Verhältnisse jener Coëfficienten überall naturgemäss sind; jedenfalls aber lassen sich, wenn man nicht in der Theorie die 
\SourceLeaf{206}
Beschränkung auf ganze Werthe der Coëfficienten weglassen will, da es dabei nicht auf absolute Werthe, sondern nur auf ihr Verhältniss unter einander ankommt, allezeit ganze Zahlen finden, die den Messungsresultaten so nahe kommen, wie man nur will.

Schliesslich wollen wir noch dem oben angeführten \textsc{Seeber}'schen Lehrsatze seine geometrische Bedeutung unterlegen. Wenn ein Parallelepipedum so beschaffen ist, dass keine seiner zwölf Kanten (unter denen je vier einander gleich sind) grösser ist, weder als eine der zwölf Diagonalen von Seitenflächen (die paarweise gleich sind), noch als eine der vier Diagonalen des Parallelepipedum: so ist der mit $\surd2$ multiplicirte Rauminhalt desselben nicht kleiner, als der Rauminhalt eines aus denselben Kanten gebildeten rechtwinklichten Parallelepipedum.
\BookRule{.28\linewidth}
